! Ό,τι έχει ``!'' μπροστά είναι σχόλιο προς τα εσένα ! Ό,τι έχει
\texttt{""} τότε είναι Πράγματα που πρέπει να ηχογραφήσεις σε ένα αρχείο
\texttt{.wav} (ή ότι θες εσύ) ! Ό,τι είναι σε \texttt{()} τότε είναι
σχόλιο αλλά για να ηχογραφήσεις κάτι. Να το δεις σαν μια ιδέα ! Ό,τι
έχει \texttt{\textasciitilde{}} κάτω από ένα αρχείο/αυτάκια (``\,``)
τότε είναι το όνομα του αρχείου που θέλω ! Όπου βλέπει
\texttt{\textless{}\textgreater{}} μπορείς να φανταστείς ότι παίζω εγώ
εκείνη την στιγμή κάποιον ήχο. Το κάνω απλά για να έχεις μια ιδέα της
ροής του κειμένου. \#\# Introduction Night

\begin{description}
\tightlist
\item[``Ο προδότης να ανοίξει τα μάτια του και οι δολοφόνοι να σηκώσουν
τα χέρια τους για να τους δει ο προδότης'']
intro\_spy\_mafia.wav
\item[``Παίζει και η τρέλα'']
intro\_madness.wav
\end{description}

\subsubsection{Night}\label{night}

``Μια νύχτα πέφτει στο Παλέρμο. Κλείστε όλοι τα μάτια σας.

\begin{description}
\tightlist
\item[Οι δολοφόνοι να ανοίξουν τα μάτια τους'']
night\_start\_mafia.wav
\item[``και να δείξουν το θύμα τους'']
night\_mafia\_target.wav
\item[``Να ξανα κλείσουν τα μάτια τους.'']
night\_mafia\_close.wav
\item[``Να ανοίξει τα μάτια του ο Σερίφης και να δείξει στο άτομο που
θέλει να μάθει κάτι.'']
night\_cop\_open.wav
\item[``Να ξανα κλείσει τα μάτια του.'']
night\_cop\_close.wav
\end{description}

\subsubsection{Day}\label{day}

\begin{description}
\tightlist
\item[``Μια νύχτα ξημερώνει στο Παλέρμο. Μπορείτε να ανοίξετε τα μάτια
σας.'']
day\_start.wav
\end{description}

\begin{description}
\tightlist
\item[``Η ψηφοφορία έχει ξεκινήσει.'']
day\_voting\_start.wav
\item[``Η ψηφοφορία έχει τελειώσει.'']
day\_voting\_end.wav
\end{description}

\subsubsection{First Day}\label{first-day}

\begin{description}
\tightlist
\item[``Το άτομο το οποίο ψηφίσατε εκτός παιχνιδιού να πάρει το κινητό
στην θέση του.'']
first\_day\_elimination.wav
\end{description}

\subsubsection{Wins}\label{wins}

\begin{description}
\tightlist
\item[``Νίκησαν οι δολοφόνοι.'']
win\_mafia.wav
\item[``Νίκησε η Τρέλα.'']
win\_madness.wav
\item[``Νίκησε το χωριό.'']
win\_village.wav
\end{description}

\subsection{Προτάσεις θανάτου/Τρόποι
θανάτου:}\label{ux3c0ux3c1ux3bfux3c4ux3acux3c3ux3b5ux3b9ux3c2-ux3b8ux3b1ux3bdux3acux3c4ux3bfux3c5ux3c4ux3c1ux3ccux3c0ux3bfux3b9-ux3b8ux3b1ux3bdux3acux3c4ux3bfux3c5}

! Κάνε Improv εδώ, αρκεί να λες κάτι από τα παρακάτω ! Τις συγκεκριμένες
προτάσεις να τις ηχογράψεις εφόσον όλες οι άλλες είναι οκ.

(Βρήκατε τον συμπολίτη σας\ldots) (Ο συμπολίτης σας βρέθηκε\ldots)
(Είδατε τον συμπολίτη σας\ldots) \textasciitilde{} death\_option.wav !
να βάζεις μετά το option έναν αριθμό

\begin{description}
\tightlist
\item[``Κανένας δεν πέθανε.'']
death\_none.wav
\end{description}
